\section{本章小结}
\label{sec:smooth-summary}

本章开头问：只看得清脚下时，一步可以走多远，走完又怎样评价？回答分成三层。

光滑性给出了依据。$L$是梯度变化的限速，由它得到下降引理~\eqref{eq:smooth-descent-lemma}：一个与目标相切的二次函数从上方盖住目标。这个上界只在$L$成立的区域内有效，与凸性互不蕴含。二阶可微时最小的$L$就是Hessian谱范数的上确界。

步长规则负责检验局部预测。固定步长靠事先知道的$L$，负梯度方向上要求$\alpha<2/L$，一般方向要重新核对长度和方向导数。回溯用真实函数值检验Armijo条件，只要方向下降就有限步终止，代价是被拒绝的试探。Wolfe条件再检验落脚点的方向导数，挡住太短的步长，并保证拟牛顿更新所需的$\langle\bm{s},\bm{y}\rangle>0$；强Wolfe条件进一步限制方向导数的绝对值，供非线性共轭梯度的下降分析使用。精确线搜索只作理论对照。

比较方法要同时固定三件事：误差标准、返回的迭代点、计算量的口径。在例~\eqref{eq:smooth-running}上，同一组运行按轮数和按读取次数排出的顺序几乎相反。凸目标通常用目标差，一般非凸目标只能用驻点残差，后者只说明一阶条件满足到什么程度。式~\eqref{eq:smooth-total-work}把被拒绝的试探和内层求解都算进总计算量。

这些规则都假设了下降方向已经给定。后面几章的问题是方向从哪里来：第~\ref{chap:first-order-methods}章只用梯度和历史位移，第~\ref{chap:second-order-methods}章用Hessian，第~\ref{chap:quasi-newton}章用梯度差拼出曲率，第~\ref{chap:trust-region}章把方向和步长放在一起决定，第~\ref{chap:saddle-points}章处理非凸目标中的鞍点。各章的收敛保证需要哪些附加条件，都在本章的标准下分别说明。
